%══════════════════════════════════════════════════════
% Identidad v5 "UAM Dato" adaptada a la tesis (clase book)
%
% Origen: ~/uam-identidad-kit/UAMDato/plantilla_dato.tex (clase article).
% Paleta, escala tipográfica (0.55 frente al deck), kicker, chips,
% cabecera/pie y componentes se copian sin cambio. Lo único que se
% adapta es el nivel de los títulos:
%   plantilla \section     -> tesis \chapter  (kicker + chip + filete)
%   plantilla \subsection  -> tesis \section
%   plantilla \subsubsection -> tesis \subsection
% Compila con lualatex (ver .latexmkrc); fontspec ya se carga en
% packages/packages.tex con las caras de Helvetica Neue por nombre PostScript.
%══════════════════════════════════════════════════════

\usepackage{titlesec,titletoc}
\usepackage{booktabs,array,multirow,colortbl,xltabular}
\usepackage{parskip}
\usepackage[most]{tcolorbox}
\tcbuselibrary{skins,breakable,raster}

% Interlineado 1.4x, el mismo del cuerpo del deck (28pt sobre 20pt).
\linespread{1.13}

%──────────────────────────────────────────────────────
% Paleta --- idéntica al deck; verificada contra el Acuerdo 06/2012.
%──────────────────────────────────────────────────────
\definecolor{uamred}{RGB}{205,3,46}        % Pantone 186 C, Azcapotzalco, numeral 5.3 p.49
\definecolor{uamblue}{RGB}{28,28,28}       % tinta de cuerpo; #000000 se reserva al emblema (numeral 2.3)
\definecolor{uamgray}{RGB}{100,100,100}
\definecolor{uamlightgray}{RGB}{244,244,244}
\definecolor{hairline}{RGB}{222,222,222}
\definecolor{darktext}{HTML}{EDEDED}
\definecolor{darkmeta}{HTML}{9A9A9A}
\definecolor{darkrule}{HTML}{3A3A3A}
\arrayrulecolor{hairline}
\AtBeginDocument{\color{uamblue}}

% Emblemas copiados del kit a images/titlepage/ para que la tesis no
% dependa de rutas fuera del proyecto.
\newcommand{\assetPie}{titlepage/logo_footer.png}
\newcommand{\assetPortada}{titlepage/logo_uam_reconstruido.pdf}

%══════════════════════════════════════════════════════
% Escala tipográfica (LetterSpace es porcentaje del em: 8 = 0.08 em)
%══════════════════════════════════════════════════════
\newcommand{\escKicker}{\bfseries\fontsize{8}{10}\selectfont
  \addfontfeature{LetterSpace=8}}
\newcommand{\escTitulo}{\bfseries\fontsize{28}{32}\selectfont}
\newcommand{\escSeccion}{\bfseries\fontsize{24}{28}\selectfont}
\newcommand{\escSubtitulo}{\fontsize{12}{16}\selectfont}
\newcommand{\escMeta}{\ttfamily\fontsize{8}{11}\selectfont}

\newcommand{\kicker}[1]{{\escKicker\color{uamred}\MakeUppercase{#1}}\par\vspace{4pt}}
\newcommand{\kickergris}[1]{{\escKicker\color{uamgray}\MakeUppercase{#1}}\par\vspace{4pt}}

%──────────────────────────────────────────────────────
% Chips circulares (la fuente se fija antes del tikzpicture para que
% `em` sea el de la letra del chip).
%──────────────────────────────────────────────────────
\newcommand{\chipgen}[3]{%
  {\bfseries\fontsize{#1}{#1}\selectfont
   \tikz[baseline=(chipnodo.base)]{%
     \node[circle, fill=#2, inner sep=0pt, minimum size=2.3em,
       text height=0.68em, text depth=0pt, align=center,
       text=white, font=\bfseries] (chipnodo) {#3};}}%
}
\newcommand{\seccionchip}[1]{\chipgen{7}{uamred}{#1}}
\newcommand{\chiplg}[1]{\chipgen{11}{uamred}{#1}}
\newcommand{\chipgris}[1]{\chipgen{7}{uamgray}{#1}}
\newcommand{\puntorojo}{\raisebox{0.22em}{\tikz\fill[uamred] (0,0) circle (2.1pt);}}
\newcommand{\puntogris}{\raisebox{0.22em}{\tikz\fill[uamgray] (0,0) circle (2.1pt);}}
\newcommand{\puntohueco}{\raisebox{0.22em}{\tikz\draw[uamgray,line width=0.5pt] (0,0) circle (2.1pt);}}

%──────────────────────────────────────────────────────
% Capítulos: kicker + chip + título en caja normal + filete, en un solo
% bloque [display] para que un salto de página no los separe. El kicker
% por omisión es "Capítulo" / "Anexo"; \bloque{...} antes de \chapter
% lo sustituye por el nombre del bloque al que pertenece.
%──────────────────────────────────────────────────────
\newcommand{\kickervacio}{}
\newcommand{\kickeractual}{}
\newcommand{\bloque}[1]{\gdef\kickeractual{#1}}
\titleformat{\chapter}[display]
  {\color{uamblue}}
  {{\escKicker\color{uamred}%
     \ifx\kickeractual\kickervacio\MakeUppercase{\chaptertitlename}%
     \else\MakeUppercase{\kickeractual}\fi}}
  {5pt}
  {\escSeccion\seccionchip{\thechapter}\hspace{0.5em}}
  [\vspace{5pt}{\color{hairline}\hrule height 0.5pt}\vspace{6pt}%
   \global\let\kickeractual\kickervacio]
% \chapter* no actualiza la marca de la cabecera: \marcasinnumero la fija
% con el propio título (titlesec se lo pasa como argumento).
\newcommand{\marcasinnumero}[1]{#1\markboth{#1}{}}
\titleformat{name=\chapter,numberless}[display]
  {\color{uamblue}}
  {}
  {0pt}
  {\escSeccion\marcasinnumero}
  [\vspace{5pt}{\color{hairline}\hrule height 0.5pt}\vspace{6pt}]
% Cada capítulo abre página nueva (book): sin espacio superior extra.
\titlespacing*{\chapter}{0pt}{0pt}{1.6ex}

\titleformat{\section}
  {\bfseries\fontsize{13}{17}\selectfont\color{uamblue}}
  {\textcolor{uamred}{\thesection}}
  {0.5em}{}
\titlespacing*{\section}{0pt}{3ex plus .5ex}{0.7ex}

\titleformat{\subsection}
  {\escKicker\color{uamgray}}
  {\thesubsection}{0.6em}{\MakeUppercase}
\titlespacing*{\subsection}{0pt}{2.4ex plus .3ex}{0.5ex}

%──────────────────────────────────────────────────────
% Cabecera y pie: filete hairline arriba y abajo, emblema y folio en
% Menlo 8 pt. \doctag es el identificador corto del documento.
%──────────────────────────────────────────────────────
\newcommand{\doctag}{ICR \textperiodcentered{} MCIE \textperiodcentered{} UAM-A}
\renewcommand{\chaptermark}[1]{\markboth{#1}{}}
\renewcommand{\sectionmark}[1]{}
\newcommand{\logopie}{\raisebox{-0.09cm}{\includegraphics[height=0.32cm]{\assetPie}}}

\pagestyle{fancy}\fancyhf{}
\fancyhead[L]{\escMeta\color{uamgray}\doctag}
\fancyhead[R]{\escMeta\color{uamgray}\nouppercase{\leftmark}}
\fancyfoot[L]{\logopie}
\fancyfoot[R]{\escMeta\color{uamgray}\thepage}
\renewcommand{\headrulewidth}{0.5pt}
\renewcommand{\footrulewidth}{0.5pt}
\renewcommand{\headrule}{{\color{hairline}\hrule width\headwidth height 0.5pt}}
\renewcommand{\footrule}{\vspace{-10pt}{\color{hairline}\hrule width\headwidth height 0.5pt}\vspace{10pt}}
\setlength{\headheight}{16pt}
% Las páginas de inicio de capítulo (estilo plain en book) llevan la misma caja.
\fancypagestyle{plain}{}

%──────────────────────────────────────────────────────
% Índice
%──────────────────────────────────────────────────────
\titlecontents{chapter}[0pt]
  {\addvspace{9pt}\bfseries\color{uamblue}}
  {\seccionchip{\thecontentslabel}\hspace{0.6em}}
  {}
  {\normalfont\color{uamgray}\;\titlerule*[0.4pc]{.}\;\contentspage}
\titlecontents{section}[1.7em]
  {\normalfont\small\color{uamblue}}
  {\textcolor{uamred}{\thecontentslabel}\hspace{0.5em}}
  {}
  {\normalfont\small\color{uamgray}\;\titlerule*[0.4pc]{.}\;\contentspage}
\titlecontents{subsection}[4.2em]
  {\normalfont\small\color{uamgray}}
  {\thecontentslabel\hspace{0.5em}}
  {}
  {\;\titlerule*[0.4pc]{.}\;\contentspage}
\titlecontents{figure}[0pt]
  {\normalfont\small\color{uamblue}}
  {\textcolor{uamred}{\thecontentslabel}\hspace{0.8em}}
  {}
  {\normalfont\small\color{uamgray}\;\titlerule*[0.4pc]{.}\;\contentspage}
\titlecontents{table}[0pt]
  {\normalfont\small\color{uamblue}}
  {\textcolor{uamred}{\thecontentslabel}\hspace{0.8em}}
  {}
  {\normalfont\small\color{uamgray}\;\titlerule*[0.4pc]{.}\;\contentspage}

\captionsetup{font=small,labelfont=bf,labelsep=period}

%══════════════════════════════════════════════════════
% Componentes del sistema (disponibles para el texto)
%   \dato / \datochico / \datohero dentro de datorow
%   \tarjeta / \tarjetasdos / \tarjetastres
%   lista filas con \filaitem{}{}
%   tablaA / tablaB con \thA / \thB / \celdaacento / \celdameta
% En la tesis la columna Y es centrada (preamble/macros.tex); en las
% tablas de este sistema usa Z, la X alineada a la izquierda.
%══════════════════════════════════════════════════════
\newcommand{\datoetiqueta}[1]{%
  {\color{uamgray}\bfseries\fontsize{8}{11}\selectfont
   \addfontfeature{LetterSpace=7}\MakeUppercase{#1}}}
\newcommand{\dato}[3][uamblue]{%
  \begin{minipage}[t]{\linewidth}\centering
    {\color{#1}\bfseries\fontsize{34}{34}\selectfont #2}\par\vspace{7pt}
    \datoetiqueta{#3}\par
  \end{minipage}%
}
\newcommand{\datochico}[3][uamblue]{%
  \begin{minipage}[t]{\linewidth}\centering
    {\color{#1}\bfseries\fontsize{23}{25}\selectfont #2}\par\vspace{6pt}
    \datoetiqueta{#3}\par
  \end{minipage}%
}
\newcommand{\datohero}[3][uamred]{%
  \par\vspace{6pt}
  \begin{minipage}[t]{\linewidth}\centering
    {\color{#1}\bfseries\fontsize{80}{80}\selectfont #2}\par\vspace{10pt}
    \datoetiqueta{#3}\par
  \end{minipage}\par\vspace{6pt}%
}
\newcolumntype{H}{!{\color{hairline}\vrule width 0.5pt}}
\newenvironment{datorow}{%
  \par\medskip\noindent{\color{hairline}\hrule height 0.5pt}\par\vspace{11pt}
}{%
  \par\vspace{9pt}{\color{hairline}\hrule height 0.5pt}\par\medskip
}

\newcommand{\tarjetaopts}{breakable}
\newtcolorbox{tarjetabase}[1][]{enhanced,
  colback=uamlightgray, colframe=uamlightgray, boxrule=0pt,
  sharp corners, left=12pt, right=12pt, top=11pt, bottom=11pt, #1}
\newcommand{\tarjeta}[3][]{%
  \begin{tarjetabase}[\tarjetaopts]
    \setlength{\parskip}{0pt}%
    \noindent\begin{minipage}[t]{\linewidth}
      \ifx\relax#1\relax\puntogris\else\chiplg{#1}\fi\hspace{7pt}%
      {\raggedright\bfseries\fontsize{10.5}{13}\selectfont\color{uamblue}#2\par}
      \vspace{5pt}
      {\raggedright\fontsize{8.5}{11.5}\selectfont\color{uamgray}#3\par}
    \end{minipage}
  \end{tarjetabase}%
}
\newcommand{\tarjetasdos}[6]{%
  \par\renewcommand{\tarjetaopts}{}%
  \begin{tcbraster}[raster columns=2, raster equal height,
    raster column skip=5mm, raster left skip=0pt, raster right skip=0pt,
    raster force size=false]
    \tarjeta[#1]{#2}{#3}
    \tarjeta[#4]{#5}{#6}
  \end{tcbraster}%
  \renewcommand{\tarjetaopts}{breakable}%
}
\newcommand{\tarjetastres}[9]{%
  \par\renewcommand{\tarjetaopts}{}%
  \begin{tcbraster}[raster columns=3, raster equal height,
    raster column skip=5mm, raster left skip=0pt, raster right skip=0pt,
    raster force size=false]
    \tarjeta[#1]{#2}{#3}
    \tarjeta[#4]{#5}{#6}
    \tarjeta[#7]{#8}{#9}
  \end{tcbraster}%
  \renewcommand{\tarjetaopts}{breakable}%
}

% Viñetas del texto corrido: el punto del sistema en lugar del cuadro que
% pone babel en español (gris, para no sumar acentos rojos).
\setlist[itemize,1]{label=\puntogris}
\setlist[itemize,2]{label=\puntohueco}

\newlist{filas}{itemize}{1}
\setlist[filas]{label=\puntohueco, leftmargin=1.5em, itemsep=9pt,
  parsep=0pt, topsep=7pt}
\newcommand{\filaitem}[2]{%
  \item {\bfseries\fontsize{11}{14}\selectfont\color{uamblue}#1}\par
        \vspace{2pt}{\fontsize{9}{12}\selectfont\color{uamgray}#2}}

\newcommand{\thA}[1]{\cellcolor{uamlightgray}%
  \raisebox{0.22\baselineskip}[0pt][0pt]{%
  {\bfseries\fontsize{8}{10}\selectfont\addfontfeature{LetterSpace=4}%
   \color{uamblue}\MakeUppercase{#1}}}}
\newcommand{\thB}[1]{%
  {\bfseries\fontsize{8}{10}\selectfont\addfontfeature{LetterSpace=4}%
   \color{uamblue}\MakeUppercase{#1}}}
\newcommand{\celdaacento}[1]{{\bfseries\color{uamred}#1}}
\makeatletter
\newcommand{\celdameta}[1]{{\ttfamily\fontsize{8}{\f@baselineskip}\selectfont
  \color{uamgray}#1}}
\makeatother
\newenvironment{tablaA}{%
  \par\medskip
  \setlength{\arrayrulewidth}{0.5pt}\arrayrulecolor{hairline}%
  \renewcommand{\arraystretch}{1.45}%
  \fontsize{10}{13.2}\selectfont
}{\par\medskip}
\newenvironment{tablaB}{%
  \par\medskip
  \setlength{\arrayrulewidth}{0.5pt}\arrayrulecolor{hairline}%
  \renewcommand{\arraystretch}{1.5}%
  \fontsize{9}{11.5}\selectfont
}{\par\medskip}
\newcommand{\reglaB}{\noalign{{\color{uamblue}\hrule height 0.6pt}}}
\newcolumntype{Z}{>{\raggedright\arraybackslash}X}
